\section{Nutzerquelle A: vollständige Runde
v1.6.3}\label{nutzerquelle-a-vollstuxe4ndige-runde-v1.6.3}

Unveränderte Quelle. Nicht alle Aussagen darin werden übernommen.
Entscheidend sind die Quellenprüfung und Korrekturen in v1.6.7.

A. Rahmen der Runde v1.6.3 Ausgangspunkt ist v1.6.2 §14 mit den vier
Prioritäten in deiner Reihenfolge: (1) Operationssatz und Kommutant, (2)
nativer Grundzustand und geladene Antwort darauf, (3) Feldwörterbuch am
Tensor, (4) räumliche Skalierung erst danach. (4) ist ausdrücklich nicht
Gegenstand. Modell unverändert: H = ΔN\_b + gΣ\_A(b\_A†P\_A +
P\_A†b\_A), P\_A = Σ\_\{i\textless j\}W\_\{A,ij\}f\_jf\_i, 64
Fermionmoden, 60 Bosonmoden, N = N\_f + 2N\_b. Kein neuer Term.

Fundament common.py (PASS, 6 Prüfgruppen): Tensor W repo-lokal, SHA
3f00a089\ldots; WWᵀ = 8·I₆₀; 60 Kanäle × 8 disjunkte Paare auf 16 Moden,
jede Mode in genau 15 Kanälen; Gewichtserhaltung q\_i+q\_j = q\_A auf
allen 480 Paaren; Casimir-Identität 8WᵀW + 𝒞\_S + 𝒞\_C − 120·I = 0 auf
allen 2016 Zweiteilchenzuständen; induzierte Bosongeneratoren X\_B =
WΛ²(X)Wᵀ/8 gaußganzzahlig mit exakter Kovarianz; gepinnter Clock (SHA
9bf99de7\ldots) mit Periode 6, Slot-Permutation {[}2,0,1,4,3{]},
Vorzeichen −1, WΛ²G\_F = G\_B·W exakt.

B. Operationssatz und Kommutant (operations\_commutant.py, PASS, 162
Guards, 0 Bareiss-Fallbacks) Frage: Welche „unsichtbare`` Freiheit
bleibt, wenn nur bestimmte Operationen verfügbar sind? Gemessen als
Dimension des Kommutanten der von der Stufe erzeugten Algebra auf dem
N-Sektor. Kleiner Kommutant = mehr unterscheidbar.

Spektralzertifikat N=3 (exakt): S = C₃C₃ᵀ annulliert durch ∏(S−r) für r
∈ \{0,7,10,12\} (Residuum 0 nnz), Multiplizitäten 64/2880/576/320 aus
rationalen Projektorspuren, tr S = 29760, tr S² = 244800.

Irreduzibel-Zerlegung (Höchstgewichtsvektoren, exakt; Gewichte in
verdoppelten Koordinaten):

Sektor Block Dim Zerlegung unter so(10)⊕\allowbreak{}su(4) Casimir (𝒞\_S, 𝒞\_C) N=2
helle Paare↔Bosonen 60 (10, 6) ×1 (36, 20) N=2 dunkle Paare 1956 (120,
10) ⊕\allowbreak{} (126, 6), je ×1 (84,36), (100,20) N=3 dunkle Tripel = ker C₃ 37888
(560, 20″) ⊕\allowbreak{} (1200, 20) ⊕\allowbreak{} (672, 4̄), je ×1 (117,63), (141,39), (165,15)
N=3 ker S (der χ-Raum) 64 (16, 4̄) ×1 (45, 15) N=3 λ=12 320 (16, 20) ×1
(45, 39) N=3 λ=10 576 (144, 4̄) ×1 (85, 15) N=3 λ=7 2880 (144, 20) ×1
(85, 39) Alle Multiplizitäten sind 1; Summenprobe über
Weyl-Dimensionsformel bestanden. Meine vorab formulierte Hypothese
stimmt in allen Blöcken.

Kommutanten-Leiter:

Stufe (verfügbare Operationen) N=2 N=3 A: nur \{X, N\_b\} (die zwei
bisher geprüften Kontrollen) 3 829 536 1 444 233 216 B: A + Clock
(Ordnung 6) 641 760 240 742 144 Cc: A + 15 su(4)-Generatoren 30 376 2
247 168 Cs: A + 45 so(10)-Generatoren 172 1 648 C: A + alle 60
Generatoren 3 7 D: B + C 3 7 volle Algebra B(H\_N) 1 1 Lesart: Die zwei
Kontrollen lassen \textasciitilde1,4·10⁹ unsichtbare Freiheitsgrade in
N=3; der Clock allein reduziert um einen Faktor 6 (seine
Eigenwertmultiplizitäten sind fast gleichverteilt, z.B.
6384/6280/6280/6384/6280/6280 auf dem Dunkelraum); erst die volle innere
Symmetrie als Operationssatz kollabiert alles auf 7 Skalare - je einen
pro (Energieblock, Irrep). so(10) ist dabei die weit wirksamere Hälfte
(1648 vs.~2,2 Mio). Das ist genau die in v1.6.2 §10.3 geforderte
Abgrenzung: Mehrfachwiederholung der zwei Kontrollen beseitigt nichts;
die Existenz der Symmetrieoperationen als Operationen ist damit nicht
bewiesen, aber ihr Effekt ist jetzt exakt beziffert.

Clock-Befund: Der Clock wirkt auf keinem irreduziblen Block als Skalar
(er verschiebt Gewichte, Zeuge z.B. {[}−2,−2,0,−2,0{]}); er liegt also
nicht in der zusammenhängenden Symmetriegruppe, sondern ist ein äußeres
Element (Slot-Permutation). D = C gilt trotzdem, weil alle
Multiplizitäten 1 sind und der Clock jeden Block in sich abbildet.

C. Nativer Grundzustand N=64 (native\_ground\_state.py, PASS, 128
Guards: 119 exakt, 9 numerisch) C.1 Lochbild (Operatoridentität,
geguardet). Mit h\_i = f\_i† ist \textbar F⟩ =
f\_0†\ldots f\_63†\textbar0⟩ das Lochvakuum und P\_A = −Q\_A†, Q\_A† =
Σ\_\{i\textless j\}W\_\{A,ij\}h\_i†h\_j†. Unter (−1)\^{}\{N\_b\}: H ≅
H\_hole = ΔN\_b + gΣ(b\_A†Q\_A† + Q\_Ab\_A). Die Teilchen-Loch-Unitäre
trägt das Vorzeichen s(μ) = (−1)\^{}\{i+j\} pro Zweilochzustand;
Negativkontrolle: ohne s(μ) ist die Komplement-Abbildung keine
Identität.

C.2 Exakte Krylov-Kette v\_n = (T†)ⁿ\textbar F⟩, T† = Σb\_A†Q\_A†:

n Einträge in v\_n ν\_n = ‖v\_n‖² ν\_n/ν\_\{n−1\} 0 1 1 - 1 480 480 480
2 108 240 439 680 916 = 4·229 3 15 252 960 575 078 400 299520/229 ≈
1307,95 4 (nicht speicherbar) 952 296 652 800 (Spurnetzwerk) ≈ 1656,0 ν₃
wurde zweifach unabhängig erhalten: durch die
15-Millionen-Einträge-Vektorrechnung und durch die
Onishi/Wick-Spurnetzwerkformel (Abschnitt E). Exakte Übereinstimmung.

C.3 Exakte Guards: T v₁ = 480 v₀; T v₂ = 916 v₁ (Eindeutigkeit des
Ein-Boson-Singuletts (60⊗Λ²64̄)\^{}G, zuvor gruppentheoretisch
vorhergesagt); Adjungiertheit ⟨v₂\textbar Tv₃⟩ = ν₃; v₁, v₂ werden von 3
so(10)- und 3 su(4)-Generatoren annulliert (Singuletts);
⟨v₁\textbar Σ\_A Q\_AQ\_A†\textbar v₁⟩ = 458·480 = 219 840 - bestätigt
die korrigierte Formel Σ\_A Q\_AQ\_A† = 480 − (15N\_h +
C\_std\^{}\{Loch\})/2 mit dem Loch-Casimir 14 der (10,6) (meine erste
Annahme „480 − 15k`` war falsch und ist durch die Rechnung ersetzt).

C.4 Schließungstest. T v₃ = α v₂ + w₂ mit α = 299520/229 und w₂ ⊥ v₂,
‖w₂‖² = 5001523200/229 ≈ 2,184·10⁷ (exakt, Pythagoras-Identität ‖w₂‖² =
‖Tv₃‖² − ν₃²/ν₂ als Guard). Defekt ‖w₂‖²/‖Tv₃‖² = 2,90·10⁻⁵. Folge: Der
Zwei-Boson-Singulettraum ist \textgreater{} 1-dimensional, die Kette ist
exakt nicht geschlossen und das Modell nicht allein durch die ν\_n
lösbar; numerisch ist die Korrektur vernachlässigbar (Kopplung
g·√(‖w₂‖²/ν₃) = 0,195g, Energieänderung in der 6. Stelle).

C.5 Exakte Ritz-Matrix (Δ=1, g herausfaktorisiert, Basis v₀\ldots v₃,
w₂): Diagonale 0, 1, 2, 3, (2 für w₂); Nebendiagonalen 4√30, 2√229,
48√29770/229; w₂ koppelt nur an v₃ mit 3√598377/11908.

C.6 Ritz-Obergrenzen E\_K(64) (rigoros, Rayleigh-Ritz):

g/Δ K=1 K=2 K=3 (+w₂) Störungstheorie 2. Ordnung −480g² 1/20 −0,704159
−0,985178 −1,094232 −1,2 1/40 −0,241620 −0,288824 −0,295357 −0,3 1/100
−0,045894 −0,047851 −0,047894 −0,048 1/400 −0,002991 −0,003000 −0,003000
−0,003 Bei g/Δ = 1/20 ist die Kette nicht konvergiert (Kopplungen ≈ 1,1Δ
\ldots{} 1,6Δ je Stufe); dafür laufen ν₅ (und mit ν₄ die Verlängerung
auf K=5). Bei ≤ 1/40 ist die Konvergenz praktisch erreicht.

C.7 Observablen bei g/Δ = 1/20 auf dem K=3-Ritz-Zustand Ω (Näherung):
⟨N\_b⟩ = 0,8421 (Hellmann-Feynman ∂E/∂Δ = 0,8421 ✓); Überlapp
\textbar⟨F\textbar Ω⟩\textbar² = 0,397; Bosonzahlverteilung p₀\ldots p₃
= 0,397 / 0,397 / 0,172 / 0,034; Shannon-Untergrenze der
Boson-Loch-Verschränkung 1,66 bit (rigorose Untergrenze, da die
reduzierte Bosondichte in N\_b blockdiagonal ist). Worker-Behauptungen
„Überlapp \textgreater{} 1/4`` und „\textgreater{} 0,811 bit``:
konsistent, aber nicht konvergiert.

C.8 Geladene Ein-Teilchen-Antwort auf Ω (Guard auf r = 0, 5, 63,
Transitivität): Entnahme f\_r\textbar Ω⟩ (nach N=63): Z\_h = 1 −
⟨N\_b⟩/32 = 0,97368; Addition f\_r†\textbar Ω⟩ (nach N=65): Z\_add =
⟨N\_b⟩/32 = 0,02632; Z\_h + Z\_add = 1 exakt. Mittlere Anregungsenergie
der Entnahme ε\_h = 0,0311Δ, der Addition ε\_add = 1,1497Δ. Kompositnorm
Z\_χ = 23⟨N\_b⟩/480 = 0,0404 (v1.6.2-N=4-Referenz 0,0301 - anderer
Zustand, nur zur Einordnung).

C.9 Sektorauswahl und Worker-Behauptung. Untergrenzen aus H ≥ (1−θ)ΔN\_b
− (15g²/2θΔ)N\_f (Casimir-Identität + quadratische Ergänzung), korrekt
als max\_θ min\_\{N\_b\}:

g/Δ Ritz E(64) schärfster Konkurrent N=64 bewiesen global? 1/400
−0,0029996 N=63: −0,0029531 ja 1/100 −0,047894 N=63: −0,04725 ja 1/40
−0,295357 N=63: −0,2953125 ja (Abstand 4,5·10⁻⁵) 1/20 −1,094232 N=66:
−1,2; N=62: −1,1625 nein Bewiesener Bereich mit diesen Schranken: g/Δ ≤
0,0251 (Bisektion). Die Worker-Behauptung „eindeutiger
N=64-Singulett-Grundzustand für 0 \textless{} \textbar g\textbar/Δ ≤
1/20`` ist damit nicht reproduziert und nicht widerlegt: Die verwendeten
Untergrenzen sind zu grob; bei 1/20 müsste E(64) \textless{} −1,2 (gegen
N ≥ 66) und \textless{} −1,1625 (gegen N=62) gezeigt werden. Ob die
konvergierte Kette dorthin kommt, entscheidet der laufende ν₅-Schritt.
Chemisches Potential: μ* = −E/64 = 0,0171Δ \textless{} Δ/50, die frühere
Aussage „μ = Δ/50 wählt den leeren Zustand`` bleibt konsistent.

D. Feldwörterbuch am Tensor (field\_dictionary.py, PASS, 29 Guards) D.1
Allgemeine Regel (bewiesen, auf allen 60 Zeilen explizit nachgerechnet):
Für Grassmann-Felder ψ\_\{I,a\} koppelt der antisymmetrische Tensor
W\_\{IJ\} genau an den symmetrischen Teil des Spinorkerns K\_\{ab\}.

Zuordnung Spinorkern K Symmetrie Kopplung an W gleichhändige Weyl, Boson
Skalar (0,0) ε\_\{αβ\} antisym. 0 (60/60 Zeilen) gleichhändige Weyl,
Boson (1,0) σ\^{}\{μν\}ε sym. ≠ 0 (Koeffizientenrang 32 je Komponente)
Dirac/Majorana, Skalar C antisym. 0 Dirac/Majorana, Pseudoskalar Cγ⁵
antisym. 0 Dirac/Majorana, Axialvektor Cγ\^{}μγ⁵ antisym. 0
Dirac/Majorana, Vektor Cγ\^{}μ sym. ≠ 0 Dirac/Majorana, Tensor
Cσ\^{}\{μν\} sym. ≠ 0 Der Vermittler kann also nur ein Vektor- oder
antisymmetrisches Tensorfeld sein, niemals
Skalar/Pseudoskalar/Axialvektor. Er trägt U(1)\_N-Ladung 2 (komplexes
Feld); Vorbehalt (kein Satz): ein masseloser geladener Vektor ist kein
konsistentes freies Eichfeld.

D.2 Chiralitätsgradierungen. Trägergraph: 64 Moden, 480 Kanten, Grad 15,
dreiecksfrei, nicht bipartit - ungerader Zyklus {[}16, 45, 0, 30, 37{]}
(Länge 5, alle Kanten geprüft; unabhängig per BFS bestätigt). Folge: Es
gibt keine globale Chiralitätsaufteilung, unter der alle 60 Bosonen
Vektoren wären. Für jede der 44 getesteten G-kovarianten Gradierungen
ist jeder Kanal rein (0 gemischte Kanäle), aber der Bosonmultiplett
zerfällt in Typen:

Gradierung Vektor-Kanäle Tensor-Kanäle Stabilisator-Dim. Kommutant der
64 Identität (keine Aufspaltung) 0 60 60 1 (Schur; direkt gegengeprüft)
Pati-Salam-Spinorsplit Γ\_pq (10 Varianten) 24 36 36 = so(6)⊕\allowbreak{}so(4)⊕\allowbreak{}su(4)
2 Farbsplit (3 Varianten) 40 20 52 = so(10)⊕\allowbreak{}(su(2)²⊕\allowbreak{}u(1)) 2 gemischt (30
Varianten) 32 28 28 4 Ein einheitlicher Lorentztyp für alle 60
Vermittler erzwingt die gleichhändige Zuordnung → nur (1,0)⊕\allowbreak{}(0,1), also
ein antisymmetrischer Tensor B\_\{μν\}. Das ist kein symmetrisches
Spin-2-Feld. Kinetische Terme: einer pro irreduziblem Block des
Stabilisators (1, 2 oder 4), keine Familienaufspaltung 4 → 1+3.

E. Spurnetzwerk-Normen (norms\_by\_traces.py, läuft: n=5) Methode:
ν\_n/(n!)² = Σ\_\{\textbar m\textbar=n\} ‖Q†\^{}m F‖²/m! über bosonische
Kohärenzzustände, Onishi-Determinante det(1 − M̄M)\^{}\{1/2\} =
exp(−½Σ\_k tr((M̄M)\^{}k)/k) und Wick-Paarung → Summe über Zyklentypen
und Permutationen von Spurnetzwerken; alle Vorfaktoren exakte Fraction.
Auswertung auf dem Φ-Tensor Φ\_\{ij,kl\} = Σ\_A(M\_A)\{ij\}(M\_A)\{kl\}
(Einträge in \{−1,0,1\}, geguardet) mit optimalem Pfad; Exaktheit über
a-priori-Schranke \textless{} 2⁵³ pro Schritt, sonst Fallback modulo
fünf Primzahlen \textless{} 2¹³ mit CRT (Eindeutigkeit über
\textbar Netzwerk\textbar{} ≤ 60ⁿ·64\^{}\{\#Zyklen\}). Validierung: ν₁ =
480, ν₂ = 439 680 (brute force im Besetzungsbild), ν₃ = 575 078 400
(identisch mit der Vektorrechnung), CRT-Route allein reproduziert ν₃.
n=4 in 22 s (112 kanonische Netzwerke), ν₄ = 952 296 652 800. n=5 läuft
seit 17 CPU-Minuten.

F. Reparaturen an den Worker-Skripten (relevant für die Bewertung der
Zahlen) Vor dem Cursor-Absturz hinterlassene Skripte hatten folgende
Fehler, alle behoben und durch Guards abgesichert: Zyklus-Slots im
Spurnetzwerk (ergab ν₂ = 466560 statt 439680); Bosonengewicht ∏m\_A!
konstant 1; T wirkte bei wiederholtem Boson doppelt (Faktor 4 statt 2 →
T nicht adjungiert zu T†, ⟨v₁\textbar Tv₂⟩ = 453120 statt 439680);
Merge-Join mit inkonsistenter Ordnung für Loch in Mode 63;
Teilchen-Loch-Guard ohne s(μ); Guard ⟨v₁\textbar ΣQQ†\textbar v₁⟩ als
Norm der Summe statt Summe der Normen; Sektorschranke min-max statt
max-min (unzulässig optimistisch) und als Stub; w₂-Norm um q² = 229²
skaliert und an falsches Kettenglied gekoppelt; sympy-Eigenwerte für
≥5×5 Radikalmatrizen (Hänger); JSON-Tupelschlüssel; Schur-Kommutant über
instabile inkrementelle SVD; ungültiger „ungerader`` 4-Zyklus-Zeuge;
komplexe Generatoren mit vermischten Real-/Imaginärteilen im
Singulett-Guard.

Nebenbefund: Das replay\_manifest.json der Vorrunde v1.6.2 war um 10:06
durch einen Lauf ohne Zugriff auf den gesperrten Documents-Tensorpfad
auf FAIL gesetzt und normal.json geleert. Aus den intakten
optimized.json (PASS-Bedingung war Byteidentität) rekonstruiert; Summen
27143/27139/4 stimmen mit dem Delivery-Manifest.

G. Einordnung gegenüber T1-T8 Tor Beitrag dieser Runde Weiter offen T1
Exakte Kommutanten-Leiter A→B→Cc→Cs→C; Clock als äußeres Element
identifiziert Welche Symmetrieoperationen der Compiler tatsächlich
bereitstellt T2 Lochbild bewiesen; Z\_h + Z\_add = 1; ε\_h, ε\_add auf
dem Ritz-Zustand Konvergierter Zustand, Renormierung, Grenzraum T3 -
räumlicher Träger (ausgeklammert) T4 Vollständige Kern-Symmetrieregel;
Nichtbipartitheit → kein einheitliches Vektor-Wörterbuch chirales Maß,
Anomalien T5 Exakte Krylov-Kette, ν₁\ldots ν₄ zweifach unabhängig
Kontinuumslimes T7 (1,0)⊕\allowbreak{}(0,1) ist B\_\{μν\}, kein Spin 2 dynamischer
Spin 2 T8 Sektorauswahl bewiesen für g/Δ ≤ 0,0251; Worker-Bereich 1/20
nicht reproduziert; μ*-Fenster Zustandsfunktional, native Präparation
Kein Tor geschlossen. Keine RH-, Faktorisierungs- oder P-vs-NP-Aussage.

H. Was noch aussteht ν₅ (läuft) → Verlängerung der Ritz-Kette auf K=5,
Neubewertung von E(64) bei 1/20 und der Sektorfrage. replay.py über alle
fünf Prüfer normal + -OO mit Byteidentität, replay\_manifest.json.
RESULTS.md / EINFACH\_ERKLAERT.md, Kopie nach universal\_room/new/ als
v1.6.3, Eintrag in universal\_room/README.md.

\begin{center}\rule{0.5\linewidth}{0.5pt}\end{center}

\section{Nutzerquelle B: Korrektur und
Zwischenbericht}\label{nutzerquelle-b-korrektur-und-zwischenbericht}

Unveränderte Quelle. Nicht alle Aussagen darin werden übernommen.
Entscheidend sind die Quellenprüfung und Korrekturen in v1.6.7.

Korrektur verifiziert: S0 jetzt \textbf{3 829 536 / 1 444 233 216}
(deckt sich mit dem v1.6.2-Originalwert und der Parallelsession), 307
Guards, PASS in beiden Modi. Hier der ausführliche Zwischenbericht.

\begin{center}\rule{0.5\linewidth}{0.5pt}\end{center}

\section[Historischer Zwischenbericht vom 15.09.2026]{Zwischenbericht: TFPT/Universalraum - fundamentale Lösung,
Stand 15.09.2026,
\textasciitilde11:45}\label{zwischenbericht-tfptuniversalraum---fundamentale-luxf6sung-stand-15.09.2026-1145}

\subsection{0. Gesamtlage: zwei Spuren}\label{gesamtlage-zwei-spuren}

Es existieren jetzt \textbf{zwei Contract-Ordner}, die sich unabhängig
auf denselben Kernpunkt zubewegt haben:

{\def\LTcaptype{none} % do not increment counter
\begin{longtable}[]{@{}
  >{\raggedright\arraybackslash}p{(\linewidth - 4\tabcolsep) * \real{0.3333}}
  >{\raggedright\arraybackslash}p{(\linewidth - 4\tabcolsep) * \real{0.3333}}
  >{\raggedright\arraybackslash}p{(\linewidth - 4\tabcolsep) * \real{0.3333}}@{}}
\toprule\noalign{}
\begin{minipage}[b]{\linewidth}\raggedright
Ordner
\end{minipage} & \begin{minipage}[b]{\linewidth}\raggedright
Ursprung
\end{minipage} & \begin{minipage}[b]{\linewidth}\raggedright
Stand
\end{minipage} \\
\midrule\noalign{}
\endhead
\bottomrule\noalign{}
\endlastfoot
\codeword{universalraum-operations-groundstate-20260915} & \textbf{diese
Session} (meine drei Arbeitspakete) & WP-A + WP-C fertig, WP-B-Sonde
läuft \\
\codeword{universalraum-native-ground-response-20260915} &
\textbf{Parallelsession} (andere Cursor-Sitzung, v1.6.4) & Grundzustand
+ geladene Antwort + Polsatz bereits konsolidiert \\
\end{longtable}
}

Die Parallelsession hat dabei einen extern eingegangenen \textbf{Polsatz
v1.6.3} reproduziert (317 Prüfbedingungen, byteidentisch) und eigene
Verschärfungen beigetragen. Beide Spuren validieren sich gegenseitig -
mit einer unten dokumentierten Abweichung, die ich inzwischen korrigiert
habe.

\textbf{Modellvertrag (unverändert in beiden Spuren):}
\(H=\Delta N_b+g\sum_A(b_A^\dagger P_A+P_A^\dagger b_A)\),
\(P_A=\sum_{i<j}W_{A,ij}f_jf_i\), \(N=N_f+2N_b\), 64 Fermionmoden, 60
Bosonmoden, \(WW^\dagger=8I_{60}\), 480 Vertizes, Prüfpunkt
\(g/\Delta=1/20\), kein \(\mu N\).

\begin{center}\rule{0.5\linewidth}{0.5pt}\end{center}

\subsection{\texorpdfstring{1. Fundament meiner Session:
\codeword{native\_source.py} (neu,
npz-frei)}{1. Fundament meiner Session:  (neu, npz-frei)}}\label{fundament-meiner-session-native_source.py-neu-npz-frei}

Problem: der gepinnte Tensor (\codeword{spinor\_tensors.npz}, SHA-256
\codeword{3f00a089…}) liegt unter \codeword{\textasciitilde{}/Documents}
- \textbf{für die Agenten-Shell macOS-gesperrt} (TCC). Die bisherigen
Verträge sind von dieser Shell nicht reproduzierbar.

Lösung: vollständige in-repo-Rekonstruktion aller Objekte aus der
Clifford/Außenalgebra-Konstruktion (die Rekonstruktion war in früheren
Runden bereits byte-exakt gegen das npz geprüft worden). \textbf{17
Guards, PASS}, darunter die volle Zwei-Teilchen-Casimiridentität
\(8W^\dagger W+4C_{\text{Spin}(10)}+4C_{\text{SU}(4)}=120\,I\) auf allen
2016 Paarzuständen. Exportiert: \codeword{W}, \codeword{J},
\codeword{C3}, Gewichte \codeword{FW/BW}, \codeword{BAR/ETA}, die 7
diskreten Symmetriegeneratoren, 60 Lie-Erzeuger, alle CAR-Helfer. Alle
drei Checker bauen nur darauf.

\begin{center}\rule{0.5\linewidth}{0.5pt}\end{center}

\subsection{2. WP-A: Der verfügbare Operationssatz (fertig, 307 exakte
Guards)}\label{wp-a-der-verfuxfcgbare-operationssatz-fertig-307-exakte-guards}

\textbf{Frage (v1.6.2 §10.3/§14):} Wie groß ist die unsichtbare Freiheit
(Kommutant) unter dem tatsächlich verfügbaren Operationssatz?

\subsubsection{Die Kommutanten-Leiter (N=3-Sektor, 45
504-dim)}\label{die-kommutanten-leiter-n3-sektor-45-504-dim}

{\def\LTcaptype{none} % do not increment counter
\begin{longtable}[]{@{}lrr@{}}
\toprule\noalign{}
Operationssatz & erzeugte Algebra & Kommutante \\
\midrule\noalign{}
\endhead
\bottomrule\noalign{}
\endlastfoot
S0: \{X, N\_b\} & 14 & \textbf{1 444 233 216} \\
S1: + Spin(10)×SU(4)-Rahmen & 743 583 744 & \textbf{7} \\
S2: + Clock-Lift & unverändert & 7 \\
S3: + Modenbesetzungen \{n\_r, m\_A\} & M₄₅₅₀₄ (2 070 614 016) &
\textbf{1} \\
S4: + geladene f-Instrumente & M₄₇₅₈₀ (N=2∪N=3) & 1 \\
\end{longtable}
}

N=2 analog: 3 829 536 → 3 → 3 → 1.

\textbf{Korrektur durch mich (Judge):} der Subagent hatte S0 als
\codeword{Σ 2m²} statt \codeword{Σ m²} gerechnet (1 452 961 792 statt 1
444 233 216) - der Guard war eine Tautologie. Jetzt korrigiert und gegen
den v1.6.2-Originalwert (\codeword{minimal\_interfaces.py}) sowie die
Parallelsession abgesichert; alle drei Quellen stimmen: \textbf{1 444
233 216}.

